\section{Changing the orders on selected pairs}\label{changing-the-orders-on-selected-pairs}

\textbf{Theorem UA.Edge (dihedral attachments).} Let
\(B=\langle b_i:i\in I\rangle\) be marked by nonidentity involutions and
have every Boolean deletion. Select a family \(\mathcal E\) of distinct
pairs \(S=\{i,j\}\) for which the product order \(q_S=|b_i b_j|\) is
finite and even. For each selected pair choose a positive multiple
\(m_S\) of \(q_S\), and attach the marked dihedral quotient \[
D_{2m_S}\longrightarrow B_S=D_{2q_S}.
\] The full compatible product of UA.SK is generated by its prescribed
marker tuples. It has the intersection property and all Boolean
deletions. Every selected pair has product order \(m_S\), and every
unselected pair retains its original product order. If \(B\) is finite,
then for every \(J\subseteq I\), \[
|H_J|=|B_J|\prod_{\substack{S\in\mathcal E\\S\subseteq J}}\frac{m_S}{q_S}.
\tag{E1}
\] If the base marking gives a regular hypertope, so does the
constructed marking, and its full type-preserving automorphism group is
\(H\).

\emph{Proof.} A group generated by two nonidentity involutions whose
product has finite order \(q_S\) is the dihedral group \(D_{2q_S}\); the
normal form alternates between powers of the product and a single
reflection. The indicated quotient exists because \(q_S\mid m_S\). It is
injective on each proper local parabolic, which is either trivial or
generated by one reflection. UA.D supplies IP in the base, and all input
deletions are part of the hypothesis. Thus UA.SK applies. The pair
supports form an antichain, so UA.SK.Local identifies every selected
pair subgroup with its full local dihedral group. UA.SK.Delete gives the
asserted deletions. The kernel on a selected pair has order \(m_S/q_S\),
and (E1) is UA.SK's order formula.

For an unselected pair, each extra coordinate of its product is a
reflection or the identity, because that coordinate's support cannot
contain both markers. If the base product has finite order \(q\), its
image under a deletion keeping just one of the two markers has order
two, so \(q\) is even. The extra coordinates therefore do not enlarge
its order. If the base product has infinite order, its lift has infinite
order by projection. This proves the assertion for every unselected
pair.

All local markers are dihedral reflections, hence involutions. Under the
final hypothesis, UA.SK.Geometry gives the geometric conclusion and
identifies the full type-preserving automorphism group. \(\square\)
